% Appendix B, part one (Chapters 1-10). Included by app_b_solutions.tex via \input.

\section{Probability Spaces}
\noindent{\small Chapter~\ref{ch:olasilik-uzaylari}.}

\begin{solution}{01-1}
Let $\F=\F_1\cap\F_2$. Since $\Omega$ is in both, $\Omega\in\F$. If $A\in\F$, then $A$ belongs to both families, and since both are closed under complements, $A^c\in\F_1\cap\F_2$. The same argument works for countable unions. (The same proof works for the intersection of arbitrarily many $\sigma$-algebras; the definition of the generated $\sigma$-algebra rests on this.)
Counterexample for the union: $\Omega=\{1,2,3\}$, $\F_1=\{\emptyset,\{1\},\{2,3\},\Omega\}$, $\F_2=\{\emptyset,\{2\},\{1,3\},\Omega\}$. Then $\{1\},\{2\}\in\F_1\cup\F_2$, but $\{1\}\cup\{2\}=\{1,2\}$ is in neither; the union is not closed under unions.
\end{solution}

\begin{solution}{01-2}
Applying the two-event formula to $A\cup B$ and $C$ and using distributivity gives
\[
\P(A\cup B\cup C)=\P(A)+\P(B)+\P(C)-\P(AB)-\P(AC)-\P(BC)+\P(ABC)
\]
($AB=A\cap B$). The general formula
\[
\P\Big(\bigcup_{i=1}^nA_i\Big)=\sum_{k=1}^n(-1)^{k+1}\sum_{i_1<\dots<i_k}\P(A_{i_1}\cap\dots\cap A_{i_k})
\]
follows by induction: for $B=\bigcup_{i\le n}A_i$, $\P(B\cup A_{n+1})=\P(B)+\P(A_{n+1})-\P\big(\bigcup_{i\le n}(A_i\cap A_{n+1})\big)$; apply the induction hypothesis to the two unions on the right and collect terms.
\emph{Matching problem:} let $A_i=\{\text{letter } i \text{ is in the right envelope}\}$. The probability that $k$ given letters are in the right place is $(n-k)!/n!$, and there are $\binom nk$ such choices; the inner sum equals $1/k!$. The probability that no letter is in the right place is
\[
1-\sum_{k=1}^n\frac{(-1)^{k+1}}{k!}=\sum_{k=0}^n\frac{(-1)^k}{k!}\xrightarrow{n\to\infty}e^{-1}\approx0.368 .
\]
Convergence is very fast: for $n=6$ the difference is below $10^{-4}$.
\end{solution}

\begin{solution}{01-3}
By De Morgan and Boole's inequality, $\P(\bigcap A_i)=1-\P(\bigcup A_i^c)\ge1-\sum\P(A_i^c)$.
\emph{Multiple testing:} for $m$ true null hypotheses let $A_i=\{\text{test } i \text{ does not reject falsely}\}$. If each test is carried out at level $\alpha/m$, then $\P(A_i^c)\le\alpha/m$, and the probability of no false rejection at all is $\ge1-\alpha$. This is the Bonferroni correction (see Exercise~\ref{exr:13-4}); it holds whatever the dependence between the tests, but it is conservative when the tests are positively dependent.
\end{solution}

\begin{solution}{01-4}
$\P(A\cap B^c)=\P(A)-\P(A\cap B)=\P(A)-\P(A)\P(B)=\P(A)\P(B^c)$, so $A\indep B^c$. Applying this result again to the pair $(B^c,A)$ gives $B^c\indep A^c$.
\end{solution}

\begin{solution}{01-5}
With prevalence $\pi$ the positive predictive value is $\mathrm{PPV}(\pi)=\frac{0.99\pi}{0.99\pi+0.05(1-\pi)}$. Then $\mathrm{PPV}>\frac12\iff0.99\pi>0.05(1-\pi)\iff\pi>\frac{0.05}{1.04}\approx0.0481$. So unless the prevalence exceeds about $4.8\%$, most positive results are false positives. For the plot:
\begin{lstlisting}[language=Python]
import numpy as np, matplotlib.pyplot as plt
pi = np.linspace(0.001, 0.5, 500)
ppv = 0.99*pi / (0.99*pi + 0.05*(1 - pi))
plt.semilogx(pi, ppv); plt.axhline(0.5, ls=":"); plt.axvline(0.05/1.04, ls=":")
plt.xlabel("prevalence"); plt.ylabel("P(D | +)"); plt.show()
\end{lstlisting}
\end{solution}

\section{Random Variables and Distribution Functions}
\noindent{\small Chapter~\ref{ch:rastgele-degiskenler}.}

\begin{solution}{02-1}
$g(x)=-2\ln x$ is strictly decreasing on $(0,1)$; its inverse is $x=e^{-y/2}$ with $|dx/dy|=\frac12e^{-y/2}$. By Theorem~\ref{thm:donusum}, $f_Y(y)=1\cdot\frac12e^{-y/2}$, $y>0$: this is $\text{Exp}(\tfrac12)=\text{Gamma}(1,\tfrac12)=\chi^2_2$.
\emph{Fisher's method:} the continuous $p$-values of $m$ independent tests are $\text{Uniform}(0,1)$ under $H_0$ (Theorem~\ref{thm:p-uniform}); hence $-2\sum_{i=1}^m\ln p_i\sim\chi^2_{2m}$, and large values of this statistic are evidence that ``at least one null hypothesis is false''.
\end{solution}

\begin{solution}{02-2}
$u=F(x)=1-e^{-x^k}\iff x=F^{-1}(u)=\big(-\ln(1-u)\big)^{1/k}$. Mean: $\E X=\int_0^\infty kx^ke^{-x^k}dx$; with $t=x^k$ this is $\int_0^\infty t^{1/k}e^{-t}dt=\Gamma(1+1/k)$. Simulation with $4\times10^5$ draws: $1.994$ for $k=0.5$ (theory $2$), $0.9031$ for $k=1.5$ ($0.9027$), $0.8934$ for $k=3$ ($0.8930$).
\end{solution}

\begin{solution}{02-3}
$\{\max(X,Y)\le z\}=\{X\le z\}\cap\{Y\le z\}$, so by independence $F_{\max}(z)=F(z)G(z)$. $\{\min(X,Y)>z\}=\{X>z\}\cap\{Y>z\}$, hence $F_{\min}(z)=1-[1-F(z)][1-G(z)]$. By induction, $\P(X_{(n)}\le x)=\prod_i\P(X_i\le x)=F(x)^n$.
\end{solution}

\begin{solution}{02-4}
$2e^{-2n(0.02)^2}\le0.05\iff n\ge\frac{\ln40}{2\cdot0.0004}\approx4611.1$; the smallest $n$ is $4612$. In general, with $\varepsilon_n=\sqrt{\ln(2/\alpha)/(2n)}$, the band
\[
L(x)=\max\{\hat F_n(x)-\varepsilon_n,0\},\qquad U(x)=\min\{\hat F_n(x)+\varepsilon_n,1\}
\]
satisfies $\P\big(L(x)\le F(x)\le U(x)\ \forall x\big)\ge1-\alpha$. This is a simultaneous, distribution-free confidence band valid for \emph{all} $x$.
\end{solution}

\section{Expectation and Variance}
\noindent{\small Chapter~\ref{ch:beklenen-deger}.}

\begin{solution}{03-1}
$X(\omega)=\int_0^\infty\I\{X(\omega)>t\}dt$. Since the integrand is non-negative, Fubini--Tonelli gives $\E X=\int_0^\infty\E\,\I\{X>t\}\,dt=\int_0^\infty\P(X>t)\,dt$. For the geometric distribution $\P(X>t)=(1-p)^{\lfloor t\rfloor}$ is a step function, and the integral is $\sum_{k=0}^\infty(1-p)^k=1/p$.
\end{solution}

\begin{solution}{03-2}
$\E[Y\mid X]=X/2$ and $\Var(Y\mid X)=X^2/12$. Tower rule: $\E Y=\E[X/2]=\frac14$. Total variance: $\Var Y=\E[X^2/12]+\Var(X/2)=\frac{1}{36}+\frac{1}{48}=\frac{7}{144}\approx0.04861$. Monte Carlo with $10^6$ draws: $\E Y\approx0.2499$, $\Var Y\approx0.04860$.
\end{solution}

\begin{solution}{03-3}
$\E(X-c)^2=\Var X+(\E X-c)^2$ is minimised at $c=\E X$, with minimum $\Var X$. $g(c)=\E|X-c|$ is convex with right derivative $g'_+(c)=\P(X\le c)-\P(X>c)=2F(c)-1$, which is $\le0$ for $c<m$ and $\ge0$ for $c\ge m$; hence $g$ is minimised at the median.
\emph{Bayes connection:} applying the same computations to the posterior distribution, the Bayes estimator is the posterior mean under squared loss and the posterior median under absolute loss (Theorem~\ref{thm:bayes-tahmin}).
\end{solution}

\begin{solution}{03-4}
$\Var Y=\rho^2+(1-\rho^2)=1$ and $\Cov(X,Y)=\rho\Var Z_1=\rho$, so $\Corr=\rho$. Since $(X,Y)$ is a linear transformation of $(Z_1,Z_2)$, it is jointly normal. $Y-\rho X=\sqrt{1-\rho^2}Z_2$ is independent of $X=Z_1$; therefore $\E[Y\mid X]=\rho X+\E[\sqrt{1-\rho^2}Z_2]=\rho X$ and $\Var(Y\mid X)=1-\rho^2$. (A special case of Theorem~\ref{thm:mvn}(4); this construction is also the standard way of generating a bivariate normal.)
\end{solution}
